% Options for packages loaded elsewhere
\PassOptionsToPackage{unicode}{hyperref}
\PassOptionsToPackage{hyphens}{url}
\PassOptionsToPackage{dvipsnames,svgnames,x11names}{xcolor}
\documentclass[
  a4paper,
]{article}
\usepackage{xcolor}
\usepackage[margin=20mm]{geometry}
\usepackage{amsmath,amssymb}
\setcounter{secnumdepth}{-\maxdimen} % remove section numbering
\usepackage{iftex}
\ifPDFTeX
  \usepackage[T1]{fontenc}
  \usepackage[utf8]{inputenc}
  \usepackage{textcomp} % provide euro and other symbols
\else % if luatex or xetex
  \usepackage{unicode-math} % this also loads fontspec
  \defaultfontfeatures{Scale=MatchLowercase}
  \defaultfontfeatures[\rmfamily]{Ligatures=TeX,Scale=1}
\fi
\usepackage{lmodern}
\ifPDFTeX\else
  % xetex/luatex font selection
  \setmainfont[]{DejaVu Serif}
  \setsansfont[]{DejaVu Sans}
  \setmonofont[]{DejaVu Sans Mono}
\fi
% Use upquote if available, for straight quotes in verbatim environments
\IfFileExists{upquote.sty}{\usepackage{upquote}}{}
\IfFileExists{microtype.sty}{% use microtype if available
  \usepackage[]{microtype}
  \UseMicrotypeSet[protrusion]{basicmath} % disable protrusion for tt fonts
}{}
\makeatletter
\@ifundefined{KOMAClassName}{% if non-KOMA class
  \IfFileExists{parskip.sty}{%
    \usepackage{parskip}
  }{% else
    \setlength{\parindent}{0pt}
    \setlength{\parskip}{6pt plus 2pt minus 1pt}}
}{% if KOMA class
  \KOMAoptions{parskip=half}}
\makeatother
\usepackage{longtable,booktabs,array}
\newcounter{none} % for unnumbered tables
\usepackage{calc} % for calculating minipage widths
% Correct order of tables after \paragraph or \subparagraph
\usepackage{etoolbox}
\makeatletter
\patchcmd\longtable{\par}{\if@noskipsec\mbox{}\fi\par}{}{}
\makeatother
% Allow footnotes in longtable head/foot
\IfFileExists{footnotehyper.sty}{\usepackage{footnotehyper}}{\usepackage{footnote}}
\makesavenoteenv{longtable}
\usepackage{graphicx}
\makeatletter
\newsavebox\pandoc@box
\newcommand*\pandocbounded[1]{% scales image to fit in text height/width
  \sbox\pandoc@box{#1}%
  \Gscale@div\@tempa{\textheight}{\dimexpr\ht\pandoc@box+\dp\pandoc@box\relax}%
  \Gscale@div\@tempb{\linewidth}{\wd\pandoc@box}%
  \ifdim\@tempb\p@<\@tempa\p@\let\@tempa\@tempb\fi% select the smaller of both
  \ifdim\@tempa\p@<\p@\scalebox{\@tempa}{\usebox\pandoc@box}%
  \else\usebox{\pandoc@box}%
  \fi%
}
% Set default figure placement to htbp
\def\fps@figure{htbp}
\makeatother
\setlength{\emergencystretch}{3em} % prevent overfull lines
\providecommand{\tightlist}{%
  \setlength{\itemsep}{0pt}\setlength{\parskip}{0pt}}
\usepackage{luatexja}
\usepackage{xurl}
\usepackage{float}
\usepackage{bookmark}
\IfFileExists{xurl.sty}{\usepackage{xurl}
\usepackage{float}}{} % add URL line breaks if available
\urlstyle{same}
\hypersetup{
  pdftitle={Two-Particle Dynamical Scattering{,} Gauge-Mediated Interactions{,} and Dynamical CHSH Bell Violation in the 3D Triangle-Loop Product (Part VI)},
  pdfauthor={Noriaki Kihara (WF System Co., Ltd.) ORCID 0009-0004-6753-4020},
  colorlinks=true,
  linkcolor={Maroon},
  filecolor={Maroon},
  citecolor={Blue},
  urlcolor={Blue},
  pdfcreator={LaTeX via pandoc}}

\title{Two-Particle Dynamical Scattering, Gauge-Mediated Interactions,
and Dynamical CHSH Bell Violation in the 3D Triangle-Loop Product (Part
VI)}
\usepackage{etoolbox}
\makeatletter
\providecommand{\subtitle}[1]{% add subtitle to \maketitle
  \apptocmd{\@title}{\par {\large #1 \par}}{}{}
}
\makeatother
\subtitle{― Two-Body Gravitational Collapse, Local Gauge Invariance
vs.~Gauge-Link Charge Splitting, Causal Spin-2 Scattering, Dynamical
Tsirelson Self-Healing (\(2\sqrt{2}\)), and Horizon Entanglement
Collapse ―}
\author{Noriaki Kihara (WF System Co., Ltd.) ORCID 0009-0004-6753-4020}
\date{November 28, 2026 Version DOI 10.5281/zenodo.24628298}

\begin{document}
\maketitle

{
\setcounter{tocdepth}{2}
\tableofcontents
}
\textbf{Author:} Noriaki Kihara (WF System Co., Ltd.)\\
\textbf{ORCID:} 0009-0004-6753-4020\\
\textbf{Date:} November 28, 2026\\
\textbf{Version DOI:} 10.5281/zenodo.24628298\\
\textbf{Concept DOI:} 10.5281/zenodo.24628297\\
\textbf{Positioning:} Part VI of the 3D triangle-loop product series
(sequel to Parts I--V, Concept DOIs \texttt{24622200},
\texttt{24622565}, \texttt{24625962}, \texttt{24626592},
\texttt{24627519}). Extending single-particle studies to spatially
separated two-particle systems and mediating wave packets on the full
dense \(H = O_8 \otimes M_8\) (\(N = 64\)) dynamics, this paper
deterministically verifies (no random numbers, zero extra assumptions,
\texttt{tri\_loop\_part6\_scattering\_chsh\_v0.py} and
\texttt{tri\_loop\_part6\_chsh\_and\_coulomb\_v0.py}): (1) two-body
gravitational collapse and local \(U(1)\) gauge redundancy
vs.~gauge-link charge splitting, (2) causal spin-2 and gauge-packet
scattering, and (3) spatially separated CHSH Bell violation and its
discontinuous collapse at the event horizon.

\begin{center}\rule{0.5\linewidth}{0.5pt}\end{center}

\subsection{Abstract}\label{abstract}

Moving beyond single-particle perturbations (Parts I--V), we place
\textbf{spatially separated two-particle excitations (sites \(A, B\))}
and mediating wave packets onto the vacuum \(H = O_8 \otimes M_8\)
(\(N = 64\)) and evolve the full nonlinear map
\(T(\tau+1) = \mathrm{Tri}(T(\tau)) / \|\mathrm{Tri}(T(\tau))\|\). We
establish five machine-precision results:

\begin{enumerate}
\def\labelenumi{\arabic{enumi}.}
\tightlist
\item
  \textbf{Two-body gravitational mass addition (sign-independent
  universal attraction) {[}SC-3{]}}: At \(\tau = 40\), the critical
  horizon collapse amplitude drops from
  \(\epsilon_c^{(1)} = \mathbf{0.273580}\) (single particle at \(A=0\))
  to \(\epsilon_c^{(2)} = \mathbf{0.269267}\) (two particles at
  \(A=0, B=1\); ratio \textbf{\texttt{0.9842}}, a 1.58\% decrease).
  Same-sign \((+,+)\) and opposite-sign \((+,-)\) pairs agree to all six
  digits, confirming universal quadratic-mass attraction.
\item
  \textbf{Local \(U(1)\) gauge redundancy vs.~gauge-link-mediated charge
  splitting {[}SC-1, SC-1b{]}}: Applying independent local \(U(1)\)
  phase rotations (\(Q = \pm 1\)) to sites \(A=0, B=1\) without a
  mediating gauge field yields an exact \textbf{\texttt{0.000e+00}}
  difference between like \((+1,+1)\) and opposite \((+1,-1)\) charges
  (exact local gauge invariance). Introducing antisymmetric \(U(1)\)
  gauge links \(G_{02} \pm G_{12}\) through a shared intermediate site
  \(C=2\) cleanly splits the nonlinear interaction strength
  \(|\Delta\|P\|_{\mathrm{int}}|\) between like \(q=(+1,+1)\) and
  opposite \(q=(+1,-1)\) pairs (\(g=0.002\): \texttt{0.2311}
  vs.~\texttt{0.0777}; \(g=0.005\): \texttt{0.0492} vs.~\texttt{0.0208};
  \(g=0.010\): \texttt{0.0933} vs.~\texttt{0.0319}).
\item
  \textbf{Causal remote two-particle spin-2 (\(\sigma_x\)) scattering
  with exact quadratic scaling {[}SC-2, SC-2b{]}}: Acting on target
  \(A=0\) from a spatially separated spin-2 (\(\sigma_x\)) source at
  \(B=1\) produces
  \textbf{\(\delta\theta(\tau=1) = \mathbf{0.0000\text{ rad}}\)} at all
  amplitudes (finite-speed causality), arrives at \(\tau=2\) with the
  exact quadratic law
  \textbf{\(\delta\theta(\tau=2) = -1.4940\,\epsilon_{\mathrm{s2}}^2\)}
  (ratio constant to 4 digits across
  \(\epsilon_{\mathrm{s2}} = 0.005 \to 0.080\), proportional to source
  energy), and accumulates up to \texttt{-2.289\ rad} by \(\tau=8\).
  Incident \(U(1)\) gauge packets \(\pm G_{01}\) exchange up to
  \texttt{+0.2907} in \(\Delta E_{\mathrm{bg}}(0)\) with sign-reversing
  phase shifts.
\item
  \textbf{Dynamical CHSH Bell violation and self-healing attractor on
  spatially separated sites (Alice=1, Bob=2) {[}SC-4b{]}}: Extracting
  the 4×4 two-qubit density matrix \(\rho_{AB}\) from the physical
  \(M_8\) doublet \(\{4,5\}\) on indices 1 and 3, every spatially
  separated site pair \((1,2), (0,1), (1,4)\) of the vacuum \(H\)
  spontaneously saturates the Tsirelson bound with purity
  \texttt{1.0000}, \(S_E = 1.0000\text{ bit}\), and
  \textbf{\(S_{\mathrm{CHSH}} = \mathbf{2.828358} \approx 2\sqrt{2}\)}
  (verified identically by both the Horodecki optimum and direct
  evaluation on four explicit unit measurement vectors). Even when
  degraded by injecting a product state \(|\uparrow_1\uparrow_2\rangle\)
  (\(S_{\mathrm{CHSH}}(0) = 2.1462, S_E = 0.291\text{ bit}\)),
  \(\mathrm{Tri}(T)\) dynamically restores entanglement to
  \textbf{\(S_{\mathrm{CHSH}} = 2.6973 \to 2.8259 \to 2.8117\)}
  (\(S_E = 0.983\text{ bit}\)) across \(\tau = 1 \to 3 \to 40\).
\item
  \textbf{Discontinuous entanglement collapse at the event horizon
  {[}SC-5b{]}}: Sweeping a center scalar mass
  \(\epsilon_{\mathrm{mass}}\) at \(A=0\) while monitoring remote sites
  \((1,2)\) at \(\tau=40\), the weak-field coexistence phase
  (\(\epsilon_{\mathrm{mass}} \le 0.27\)) preserves
  \textbf{\(S_{\mathrm{CHSH}} = \mathbf{2.8284}\)}
  (\texttt{purity\ =\ 1.0000}), whereas crossing the horizon threshold
  \(\epsilon_c \approx 0.2736\) (\(\epsilon_{\mathrm{mass}} \ge 0.28\))
  causes a discontinuous drop to
  \textbf{\(S_{\mathrm{CHSH}} = \mathbf{0.0000}\)}
  (\texttt{purity\ =\ 0.2500}, maximally mixed \(\mathbb{I}/4\)).
\end{enumerate}

\begin{center}\rule{0.5\linewidth}{0.5pt}\end{center}

\subsection{1. Introduction: From Single Particles to Two-Body
Scattering and Bell
Nonlocality}\label{introduction-from-single-particles-to-two-body-scattering-and-bell-nonlocality}

Parts I--V {[}P1--P5{]} established that the triangle-loop product
\(\mathrm{Tri}(T)_{ijk} = \sum_{p,q,r} T_{ipq} T_{qjr} T_{rpk}\) on
\(H = O_8 \otimes M_8\) spontaneously generates one-generation Standard
Model gauge symmetries, electroweak symmetry breaking, multi-shell
gravity, and spin-2 time dilation \(g_{00}(d)\). The next essential
question is how \textbf{spatially separated two-particle systems
interact, scatter, and sustain or lose nonlocal CHSH Bell correlations
(\(S > 2\))} under the exact nonlinear dynamics of \(\mathrm{Tri}(T)\).
Each claim below is tagged as \textbf{{[}Known{]}}, \textbf{{[}Part I--V
Established{]}}, \textbf{{[}Part VI Verified{]}}, or
\textbf{{[}Hypothesis / Unverified{]}}.

\begin{center}\rule{0.5\linewidth}{0.5pt}\end{center}

\subsection{2. Two-Body Gravitational Collapse, Local Gauge Invariance,
and Gauge-Link Charge Splitting {[}Part VI Verified SC-1, SC-1b,
SC-3{]}}\label{two-body-gravitational-collapse-local-gauge-invariance-and-gauge-link-charge-splitting-part-vi-verified-sc-1-sc-1b-sc-3}

\subsubsection{\texorpdfstring{2.1 Two-Particle Horizon Threshold
\(\epsilon_c^{(2)}\) and Universal Attraction
(SC-3)}{2.1 Two-Particle Horizon Threshold \textbackslash epsilon\_c\^{}\{(2)\} and Universal Attraction (SC-3)}}\label{two-particle-horizon-threshold-epsilon_c2-and-universal-attraction-sc-3}

Bisecting the \(\tau=40\) horizon collapse threshold on
\(H = O_8 \otimes M_8\) (\(N=64\)) for scalar particles \(P_0, P_1\) at
sites \(A=0, B=1\): - Single particle (\(A=0\)):
\(\epsilon_c^{(1)} = (0.273579, 0.273580)\) (midpoint
\textbf{\texttt{0.273580}}) - Two particles, same sign (\(P_0 + P_1\)):
\(\epsilon_c^{(2,+)} = (0.269267, 0.269268)\) (midpoint
\textbf{\texttt{0.269267}}, ratio \textbf{\texttt{0.9842}}) - Two
particles, opposite sign (\(P_0 - P_1\)):
\(\epsilon_c^{(2,-)} = (0.269267, 0.269268)\) (midpoint
\textbf{\texttt{0.269267}}, ratio \textbf{\texttt{0.9842}})

Two spatially separated particles lower the collapse threshold by 1.58\%
with exact six-digit sign independence, reflecting the \(\mathbb{Z}_2\)
selection rule of \(M_8\) that couples scalar gravity strictly through
squared mass \(\propto \epsilon^2\).

\subsubsection{\texorpdfstring{2.2 Local \(U(1)\) Gauge Redundancy
(SC-1) vs.~Gauge-Link Charge Splitting (SC-1b, Fig.
1(b))}{2.2 Local U(1) Gauge Redundancy (SC-1) vs.~Gauge-Link Charge Splitting (SC-1b, Fig. 1(b))}}\label{local-u1-gauge-redundancy-sc-1-vs.-gauge-link-charge-splitting-sc-1b-fig.-1b}

\begin{enumerate}
\def\labelenumi{\arabic{enumi}.}
\tightlist
\item
  \textbf{No-gauge control (SC-1)}: Independent local \(U(1)\) rotations
  (\(Q = \pm 1\)) on the \(\{4,5\}\) doublet at \(A=0, B=1\) yield zero
  difference (\texttt{0.000e+00}) between \(Q=(+1,+1)\) and
  \(Q=(+1,-1)\) across \(\epsilon = 0.05, 0.12, 0.18\), confirming exact
  local \(U(1)\) gauge invariance.
\item
  \textbf{Gauge-link mediation (SC-1b)}: Coupling \(A=0\) and \(B=1\)
  through antisymmetric \(U(1)\) gauge links \(G_{02} \pm G_{12}\) at
  shared site \(C=2\) (\(\epsilon = 0.08, \tau = 20\)) cleanly splits
  \(\Delta\|P\|_{\mathrm{int}} = \|P_{AB}\| - \|P_A\| - \|P_B\| + \|P_{\mathrm{vac}}\|\)
  between like and opposite charges:
\end{enumerate}

{\def\LTcaptype{none} % do not increment counter
\begin{longtable}[]{@{}
  >{\centering\arraybackslash}p{(\linewidth - 6\tabcolsep) * \real{0.2500}}
  >{\centering\arraybackslash}p{(\linewidth - 6\tabcolsep) * \real{0.2500}}
  >{\centering\arraybackslash}p{(\linewidth - 6\tabcolsep) * \real{0.2500}}
  >{\centering\arraybackslash}p{(\linewidth - 6\tabcolsep) * \real{0.2500}}@{}}
\toprule\noalign{}
\begin{minipage}[b]{\linewidth}\centering
Gauge coupling \(g\)
\end{minipage} & \begin{minipage}[b]{\linewidth}\centering
Like \(q=(+1,+1)\) \(\Delta\|P\|_{\mathrm{int}}\)
\end{minipage} & \begin{minipage}[b]{\linewidth}\centering
Opposite \(q=(+1,-1)\) \(\Delta\|P\|_{\mathrm{int}}\)
\end{minipage} & \begin{minipage}[b]{\linewidth}\centering
Like vs.~Opposite \(\|P\|(20)\)
\end{minipage} \\
\midrule\noalign{}
\endhead
\bottomrule\noalign{}
\endlastfoot
\texttt{0.002} & \textbf{\texttt{-2.3111e-01}} &
\textbf{\texttt{-7.7669e-02}} & \texttt{0.073727}
vs.~\texttt{0.227169} \\
\texttt{0.005} & \textbf{\texttt{-4.9178e-02}} &
\textbf{\texttt{-2.0825e-02}} & \texttt{0.085754}
vs.~\texttt{0.114107} \\
\texttt{0.010} & \textbf{\texttt{-9.3327e-02}} &
\textbf{\texttt{-3.1906e-02}} & \texttt{0.167565}
vs.~\texttt{0.228986} \\
\texttt{0.020} & \textbf{\texttt{-4.8884e-03}} &
\textbf{\texttt{-1.0441e-02}} & \texttt{0.107021}
vs.~\texttt{0.101469} \\
\end{longtable}
}

\begin{center}\rule{0.5\linewidth}{0.5pt}\end{center}

\subsection{3. Dynamical Scattering and Phase/Energy Exchange: Causal
Spin-2 and Gauge-Packet Scattering {[}Part VI Verified SC-2,
SC-2b{]}}\label{dynamical-scattering-and-phaseenergy-exchange-causal-spin-2-and-gauge-packet-scattering-part-vi-verified-sc-2-sc-2b}

For target particle \(A=0\) (\(\epsilon_t = 0.08\)) under remote spin-2
(\(\sigma_x\)) excitation at \(B=1\) or incident \(U(1)\) gauge packets
\(\pm G_{01}\) (Fig. 1(a)):

\begin{enumerate}
\def\labelenumi{\arabic{enumi}.}
\tightlist
\item
  \textbf{Remote spin-2 (\(\sigma_x\)) causality and exact quadratic
  scaling (SC-2b)}:

  \begin{itemize}
  \tightlist
  \item
    \(\tau = 1\):
    \textbf{\(\delta\theta(\tau=1) = \mathbf{+0.0000\text{ rad}}\)}
    across all \(\epsilon_{\mathrm{s2}} = 0.005 \dots 0.080\)
    (finite-speed causality).
  \item
    \(\tau = 2\) (arrival front):
    \(\delta\theta(\tau=2) = -3.735\times 10^{-5},\; -1.494\times 10^{-4},\; -5.975\times 10^{-4},\; -2.389\times 10^{-3},\; -9.532\times 10^{-3}\,\mathrm{rad}\),
    with
    \textbf{\(\delta\theta(\tau=2)/\epsilon_{\mathrm{s2}}^2 = \mathbf{-1.4940,\; -1.4940,\; -1.4937,\; -1.4929,\; -1.4893}\)}
    (4-digit constant, linear in source energy
    \(M \propto \epsilon_{\mathrm{s2}}^2\)).
  \item
    \(\tau = 8\): accumulates to \textbf{\texttt{-2.2890\ rad}} at
    \(\epsilon_{\mathrm{s2}} = 0.08\).
  \end{itemize}
\item
  \textbf{\(U(1)\) gauge-packet scattering \(\pm G_{01}\) (SC-2)}:

  \begin{itemize}
  \tightlist
  \item
    Target background energy transfer \(\Delta E_{\mathrm{bg}}(0, +G)\)
    grows from \texttt{+1.599e-03} (\(\tau=1\)) to
    \textbf{\texttt{+2.907e-01}} (\(\tau=10\)).
  \item
    Packet signs \(+G\) and \(-G\) invert the initial phase shift
    (\texttt{+1.038e-11} vs.~\texttt{-1.221e-11} at \(\tau=1\)) and
    trigger asymmetric \(\pi\text{ rad}\) resonant phase flips at
    \(\tau = 8\text{–}15\).
  \end{itemize}
\end{enumerate}

\begin{figure}[H]
\centering
\includegraphics[width=0.96\linewidth,height=\textheight,keepaspectratio,alt={Part VI Verification (1): (a) Remote two-particle spin-2 (\textbackslash sigma\_x) dynamical scattering showing causality (\textbackslash delta\textbackslash theta(\textbackslash tau=1)=0) and exact quadratic arrival scaling \textbar\textbackslash delta\textbackslash theta(\textbackslash tau=2)\textbar{} = 1.494\textbackslash,\textbackslash epsilon\_\{\textbackslash mathrm\{s2\}\}\^{}2; (b) Like q=(+1,+1) vs.~opposite q=(+1,-1) charge splitting via U(1) gauge links contrasted with the no-gauge control.}]{figs_en/fig_p6_two_particle_scattering_en.png}
\caption{Part VI Verification (1): (a) Remote two-particle spin-2
(\(\sigma_x\)) dynamical scattering showing causality
(\(\delta\theta(\tau=1)=0\)) and exact quadratic arrival scaling
\(|\delta\theta(\tau=2)| = 1.494\,\epsilon_{\mathrm{s2}}^2\); (b) Like
\(q=(+1,+1)\) vs.~opposite \(q=(+1,-1)\) charge splitting via \(U(1)\)
gauge links contrasted with the no-gauge control.}
\end{figure}

\begin{center}\rule{0.5\linewidth}{0.5pt}\end{center}

\subsection{4. Spatially Separated CHSH Bell Violation and Horizon
Entanglement Collapse {[}Part VI Verified SC-4b,
SC-5b{]}}\label{spatially-separated-chsh-bell-violation-and-horizon-entanglement-collapse-part-vi-verified-sc-4b-sc-5b}

\subsubsection{\texorpdfstring{4.1 Two-Qubit Density Matrix
\(\rho_{AB}\) and Vacuum Tsirelson Saturation
(\(2\sqrt{2}\))}{4.1 Two-Qubit Density Matrix \textbackslash rho\_\{AB\} and Vacuum Tsirelson Saturation (2\textbackslash sqrt\{2\})}}\label{two-qubit-density-matrix-rho_ab-and-vacuum-tsirelson-saturation-2sqrt2}

Tracing out the intermediate environment \((M, m)\) from the physical
\(M_8\) doublet \(\{4,5\}\) on index 1 (Alice = site \(A\)) and index 3
(Bob = site \(B\)) yields a positive-semidefinite 4×4 density matrix
\(\rho_{AB} = \mathrm{Tr}_{\mathrm{Env}}(|\Psi\rangle\langle\Psi|)\).
Across every spatially separated vacuum pair
\((A,B) = (1,2), (0,1), (1,4)\), the Pauli correlation tensor takes the
\(Y\)-rotated singlet form (\(\alpha \approx 57.7^\circ\))
\[T_{ij} = \mathrm{Tr}(\rho_{AB}(\sigma_i\otimes\sigma_j)) = \begin{pmatrix} -0.5344 & 0 & +0.8452 \\ 0 & -1.0000 & 0 \\ -0.8452 & 0 & -0.5344 \end{pmatrix},\]
achieving purity \textbf{\texttt{1.0000}},
\(S_E = \mathbf{1.0000\text{ bit}}\), and
\textbf{\(S_{\max} = S_{\mathrm{direct}} = \mathbf{2.828358} \approx 2\sqrt{2}\)}
(verified identically by the Horodecki optimum and direct trace
evaluation on four explicit unit measurement vectors).

\subsubsection{4.2 Dynamical Entanglement Generation and Bell-State
Preservation (SC-4b, Fig.
2(a))}\label{dynamical-entanglement-generation-and-bell-state-preservation-sc-4b-fig.-2a}

Injecting a product state \(|\uparrow_1\uparrow_2\rangle\) or Bell
states \(|\Phi^+\rangle, |\Psi^-\rangle\) (\(\epsilon = 0.10\)) onto
\((A=1, B=2)\) and evolving \(\tau = 0 \dots 40\):

{\def\LTcaptype{none} % do not increment counter
\begin{longtable}[]{@{}
  >{\raggedright\arraybackslash}p{(\linewidth - 8\tabcolsep) * \real{0.1667}}
  >{\centering\arraybackslash}p{(\linewidth - 8\tabcolsep) * \real{0.2083}}
  >{\centering\arraybackslash}p{(\linewidth - 8\tabcolsep) * \real{0.2083}}
  >{\centering\arraybackslash}p{(\linewidth - 8\tabcolsep) * \real{0.2083}}
  >{\centering\arraybackslash}p{(\linewidth - 8\tabcolsep) * \real{0.2083}}@{}}
\toprule\noalign{}
\begin{minipage}[b]{\linewidth}\raggedright
Initial injected state
\end{minipage} & \begin{minipage}[b]{\linewidth}\centering
\(\tau = 0\) \(S_{\mathrm{CHSH}}\) (\(S_E\))
\end{minipage} & \begin{minipage}[b]{\linewidth}\centering
\(\tau = 1\)
\end{minipage} & \begin{minipage}[b]{\linewidth}\centering
\(\tau = 3\)
\end{minipage} & \begin{minipage}[b]{\linewidth}\centering
\(\tau = 40\) \(S_{\mathrm{CHSH}}\) (\(S_E\))
\end{minipage} \\
\midrule\noalign{}
\endhead
\bottomrule\noalign{}
\endlastfoot
\textbf{Product \(|\uparrow_1\uparrow_2\rangle\)} & \texttt{2.1462}
(\texttt{0.291\ bit}) & \texttt{2.6973} (\texttt{0.866\ bit}) &
\textbf{\texttt{2.8259}} (\texttt{0.998\ bit}) &
\textbf{\texttt{2.8117}} (\texttt{0.983\ bit}) \\
\textbf{Bell \(|\Phi^+\rangle\)} & \texttt{2.7931} (\texttt{1.000\ bit})
& \texttt{2.8274} (\texttt{1.000\ bit}) & \texttt{2.8274}
(\texttt{1.000\ bit}) & \textbf{\texttt{2.8274}}
(\texttt{1.000\ bit}) \\
\textbf{Bell \(|\Psi^-\rangle\)} & \texttt{2.8273} (\texttt{1.000\ bit})
& \texttt{2.8275} (\texttt{1.000\ bit}) & \texttt{2.8281}
(\texttt{1.000\ bit}) & \textbf{\texttt{2.8282}}
(\texttt{1.000\ bit}) \\
\end{longtable}
}

Rather than degrading quantum correlations, the nonlinear
\(\mathrm{Tri}(T)\) dynamics acts as a \textbf{self-healing entanglement
attractor}, lifting degraded product perturbations back to the Tsirelson
bound within 3 steps.

\subsubsection{4.3 Discontinuous Entanglement Collapse at the Event
Horizon (SC-5b, Fig.
2(b))}\label{discontinuous-entanglement-collapse-at-the-event-horizon-sc-5b-fig.-2b}

Sweeping center scalar mass \(\epsilon_{\mathrm{mass}}\) at \(A=0\) and
measuring remote sites \((1,2)\) at \(\tau=40\): - \textbf{Weak-field
coexistence
(\(\epsilon_{\mathrm{mass}} = 0.00 \dots 0.27 < \epsilon_c\))}:
\texttt{purity\ =\ 1.0000}, \(S_E = 1.0000\text{ bit}\),
\textbf{\(S_{\mathrm{CHSH}} = \mathbf{2.8284}\)} (\texttt{VIOLATED}). -
\textbf{Horizon-collapsed phase
(\(\epsilon_{\mathrm{mass}} = 0.28, 0.35 > \epsilon_c\))}: simultaneous
with \(\|P\|(40) = 0.499976\), the remote pair collapses discontinuously
to \textbf{\texttt{purity\ =\ 0.2500}} (maximally mixed
\(\mathbb{I}/4\)) and \textbf{\(S_{\mathrm{CHSH}} = \mathbf{0.0000}\)}
(\texttt{CLASSICAL}).

Thus the geometric horizon threshold \(\epsilon_c \approx 0.2736\)
coincides exactly with a \textbf{discontinuous entanglement-severing
transition} (\(2\sqrt{2} \to 0\)).

\begin{figure}[H]
\centering
\includegraphics[width=0.96\linewidth,height=\textheight,keepaspectratio,alt={Part VI Verification (2): (a) Dynamical CHSH generation from a product state and Bell-state preservation on spatially separated sites (Alice=1, Bob=2); (b) Discontinuous CHSH collapse (2\textbackslash sqrt\{2\}\textbackslash to 0, purity 1.0\textbackslash to 0.25) at the event horizon threshold \textbackslash epsilon\_c \textbackslash approx 0.2736.}]{figs_en/fig_p6_chsh_and_horizon_transition_en.png}
\caption{Part VI Verification (2): (a) Dynamical CHSH generation from a
product state and Bell-state preservation on spatially separated sites
(Alice=1, Bob=2); (b) Discontinuous CHSH collapse (\(2\sqrt{2}\to 0\),
purity \(1.0\to 0.25\)) at the event horizon threshold
\(\epsilon_c \approx 0.2736\).}
\end{figure}

\begin{center}\rule{0.5\linewidth}{0.5pt}\end{center}

\subsection{5. Open Problems and
Reproducibility}\label{open-problems-and-reproducibility}

\begin{itemize}
\tightlist
\item
  \textbf{Open problems (deferred to Part VII+)}: (1) distance-dependent
  two-particle wave-packet collision cross-sections on
  \(O_8^{\otimes K}\;(K\ge 3)\); (2) three-generation structure via the
  exceptional Jordan algebra \(J_3(\mathbb{O})\), precision fermion mass
  ratios, and first-principles derivation of \(\alpha\) and \(\alpha_s\)
  (planned Part VII).
\item
  \textbf{Reproducibility files}:
  \texttt{tri\_loop\_part6\_scattering\_chsh\_v0.py} (SC-1--SC-5),
  \texttt{results/tri\_loop\_part6\_scattering\_chsh\_v0.txt},
  \texttt{tri\_loop\_part6\_chsh\_and\_coulomb\_v0.py} (SC-1b--SC-5b),
  \texttt{results/tri\_loop\_part6\_chsh\_and\_coulomb\_v0.txt},
  \texttt{plot\_part6\_figures\_v0.py}, and
  \texttt{results/fig\_p6\_*\_\{ja,en\}.png}.
\end{itemize}

\begin{center}\rule{0.5\linewidth}{0.5pt}\end{center}

\subsection{References}\label{references}

\subsubsection{A. Self-Citations}\label{a.-self-citations}

\begin{itemize}
\tightlist
\item
  \textbf{{[}P1{]}} N. Kihara, \emph{Spontaneous Emergence of Quantum,
  Gauge, Gravitational, and Fermionic Structures from the 3D
  Triangle-Loop Product (Integrated Part I)}, Zenodo (2026-11-28),
  Concept DOI: \texttt{10.5281/zenodo.24622200}.
\item
  \textbf{{[}P2{]}} N. Kihara, \emph{Unified Connection Hypothesis for
  Gravity, Electromagnetism, QCD, and Dirac Fermions in the 3D
  Triangle-Loop Product (Part II)}, Zenodo (2026-11-28), Concept DOI:
  \texttt{10.5281/zenodo.24622565}.
\item
  \textbf{{[}P3{]}} N. Kihara, \emph{Spacetime and Internal-Space
  Readout Connection Hypothesis in the 3D Triangle-Loop Product (Part
  III)}, Zenodo (2026-11-28), Concept DOI:
  \texttt{10.5281/zenodo.24625962}.
\item
  \textbf{{[}P4{]}} N. Kihara, \emph{Electroweak Symmetry Breaking and
  Maximal Parity Violation in the 3D Triangle-Loop Product (Part IV)},
  Zenodo (2026-11-28), Concept DOI: \texttt{10.5281/zenodo.24626592}.
\item
  \textbf{{[}P5{]}} N. Kihara, \emph{Multi-Shell Large-Lattice Gravity,
  Causal Propagation, and Spin-2 Unlock of \(g_{00}(d)\) in the 3D
  Triangle-Loop Product (Part V)}, Zenodo (2026-11-28), Concept DOI:
  \texttt{10.5281/zenodo.24627519}.
\end{itemize}

\subsubsection{B. External Primary
Literature}\label{b.-external-primary-literature}

\begin{itemize}
\tightlist
\item
  \textbf{{[}E1{]}} J. F. Clauser, M. A. Horne, A. Shimony, and R. A.
  Holt, \emph{Proposed Experiment to Test Local Hidden-Variable
  Theories}, Phys. Rev.~Lett. \textbf{23}, 880--884 (1969).
\item
  \textbf{{[}E2{]}} R. Horodecki, P. Horodecki, and M. Horodecki,
  \emph{Violating Bell inequality by mixed spin-1/2 states: necessary
  and sufficient condition}, Phys. Lett. A \textbf{200}, 340--344
  (1995).
\item
  \textbf{{[}E3{]}} B. S. Cirel'son (Tsirelson), \emph{Quantum
  generalizations of Bell's inequality}, Lett. Math. Phys. \textbf{4},
  93--100 (1980).
\end{itemize}

\end{document}
